(i) If $X = (x_{ij})$ commutes with every $E_{kl}$, then $E_{kk} X E_{ll} = X E_{kk} E_{ll}$ gives $x_{kl} = 0$ for $k \ne l$, and $E_{kl} X = X E_{kl}$ gives $x_{kk} = x_{ll}$. If it also commutes with $a 1_m$ for every $a$, then $x_{11} \in \Zc(A)$. Conversely such scalar matrices are central.

(ii) For a two-sided ideal $J$ of $\Mat_m(A)$ let $I$ be the set of entries of elements of $J$. Since $E_{1i} X E_{j1} = x_{ij} E_{11}$, the element $a E_{11}$ lies in $J$ for every $a \in I$, and then $a E_{kl} = E_{k1}(a E_{11})E_{1l} \in J$. So $\Mat_m(I) \subseteq J \subseteq \Mat_m(I)$, and $I$ is an ideal because $(b a c) E_{11} = (b E_{11})(a E_{11})(c E_{11})$. Conversely $\Mat_m(I)$ is an ideal.

(iii) For a left $\Mat_m(A)$-module $N$, $E_{11} N$ is a left $A$-module via $a \cdot n = (a E_{11}) n$, and $n \mapsto (E_{1i} n)_i$ is an isomorphism $N \to (E_{11}N)^{m}$ with inverse $(n_i) \mapsto \sum_i E_{i1} n_i$, because $\sum_i E_{i1}E_{1i} = 1$. For a left $A$-module $M$, $E_{11} M^{m}$ is the first coordinate, isomorphic to $M$. Both constructions are functorial and the isomorphisms are natural. Throughout, $E_{kl}$ are the matrix units of $\Mat_m(A)$ and $A \cong E_{11}\Mat_m(A)E_{11}$, by \claimref{std-algebra.matrix-units}.
